\ifdefined\mainfile\else
\documentclass{article}
\usepackage{graphicx}
\usepackage{float}
\usepackage{amsmath}
\usepackage{amssymb}
\usepackage[a4paper, left=2.5cm, right=2.5cm, top=2.5cm, bottom=2.5cm]{geometry}
\usepackage[colorlinks=true,linkcolor=black,citecolor=blue,urlcolor=blue]{hyperref}
\begin{document}
\fi

\subsection{Hardware design}
\label{sec:hardware}

The brain of the bin is an \textbf{ESP8266 NodeMCU} module - it has Wifi-capabilities, which allows us to run the intended program. It was not, the ideal candidate: the ESP8266 only exposes a handful of usable GPIO pins and several of those have boot-time roles, which made the pin layout one of the recurring headaches of the project. However we manged to keep the entire system limited to a single ESP-unit, making the system cheaper to produce and sell.
\\ \\
Around the ESP we built up the rest of the hardware listed in Table~\ref{tab:components}. The selection satisfies the course's component requirements -- one analog sensor (LM35), one digital sensor (HC-SR04), one form of actuation (the steppers) and an on-device display (the OLED).

\begin{table}[H]
    \centering
    \begin{tabular}{l|l|l}
        \textbf{Component} & \textbf{Role} & \textbf{Type} \\
        \hline
        ESP8266 NodeMCU                 & WiFi MCU                & Microcontroller \\
        HC-SR04 ultrasonic sensor       & Fill level measurement  & \textbf{Digital} sensor \\
        LM35 temperature sensor         & Temperature measurement & \textbf{Analog} sensor \\
        28BYJ-48 stepper motor (x2)     & Bag closure mechanism   & \textbf{Actuator} \\
        ULN2003 driver board (x2)       & Stepper coil driver     & Driver \\
        SSD1306 OLED display (I\textsuperscript{2}C, 0.96\textquotedbl) & Local status display & Display \\
        Push button                     & User reset              & Input \\
    \end{tabular}
    \caption{Components used in the Smart Trash Bin}
    \label{tab:components}
\end{table}
\noindent
Power is supplied from a battery and convertet down to 5\,V. The ESP8266 runs from its onboard 3.3\,V regulator, while the stepper drivers and the OLED share the 5\,V rail. The HC-SR04 outputs 5\,V on its Echo line, so it is divided down with a 1\,k$\Omega$/2\,k$\Omega$ resistor pair before reaching a 3.3\,V-only GPIO. The full pinout is listed in Table~\ref{tab:pinout}, and a larger circuit diagram is reproduced in the appendix.

\begin{table}[H]
    \centering
    \begin{tabular}{l|l|l}
        \textbf{Component} & \textbf{ESP8266 pin} & \textbf{Note} \\
        \hline
        HC-SR04 Trig     & D5  & digital output \\
        HC-SR04 Echo     & D6  & digital input (1k/2k divider) \\
        ULN2003 IN1      & D8  & stepper coil 1 \\
        ULN2003 IN2      & D7  & stepper coil 2 \\
        ULN2003 IN3      & D2  & stepper coil 3 (shared with I\textsuperscript{2}C SDA) \\
        ULN2003 IN4      & D0  & stepper coil 4 (moved from D1 to free SCL) \\
        SSD1306 SDA      & D2  & I\textsuperscript{2}C data \\
        SSD1306 SCL      & D1  & I\textsuperscript{2}C clock \\
        LM35 Vout        & A0  & analog input, 10-bit ADC \\
        Reset button     & D3  & digital input, internal pull-up \\
    \end{tabular}
    \caption{Pinout for the ESP8266 module}
    \label{tab:pinout}
\end{table}
\noindent
\noindent To check whether the bin could run without mains power we made a small power budget in MATLAB. The script sums the current drawn by the ESP8266, the OLED, the two sensors and the stepper driver, weights each by how long the bin spends in every state over a day, and includes the loss in the buck converter. Comparing a 9\,V alkaline battery with a pair of 18650 Li-ion cells, it shows that the alkaline option lasts well under a day with the WiFi always on, while the Li-ion pack lasts considerably longer, and that a modem-sleep or light-sleep mode would extend this further. Deep sleep was left out on purpose, since on the ESP8266 it resets the chip on wake and would lose the state machine.
\begin{align*}
Total_{Power} &= V_{5\text{V}}\ I_{\text{avg}} & I_{\text{avg}} &= \frac{1}{T}\sum_{s} I_{s}\, t_{s} \\
I_{\text{batt}} &= \frac{Total_{Power}}{\eta\, V_{\text{batt}}} + I_{q} & t_{\text{life}} &= \frac{Q_{\text{batt}}}{I_{\text{batt}}}
\end{align*}
\noindent where $I_{s}$ and $t_{s}$ are the average current and the time spent in state $s$ over one day $T$, $V_{5\text{V}}$ is the rail voltage, $\eta$ the buck-converter efficiency, $I_{q}$ the buck quiescent current, and $Q_{\text{batt}}$ and $V_{\text{batt}}$ the battery charge and nominal voltage.
\noindent The graphs for the power calculations, can be found on the poster.
\ifdefined\mainfile\else
\end{document}
\fi
